% STATUS: DRAFT
\chapter{Inequivalent representations}\label{ch:inequiv}

\begin{physmotiv}
Two ferromagnets, one magnetized up and one down, are described by the \emph{same} algebra of local observables but not by the same Hilbert space: no finite number of local operations turns one into the other. Phases of matter, superselection sectors, spontaneous symmetry breaking and thermal states in infinite volume are all instances of one mathematical fact: a \cstar-algebra has many inequivalent representations, and GNS (\cref{ch:gns}) shows that they correspond to different states. This chapter makes that precise and gives practical criteria.
\end{physmotiv}

\section{Equivalence, disjointness, quasi-equivalence}

Two representations $\pi_1,\pi_2$ of $\A$ on $\HH_1,\HH_2$ are \emph{unitarily equivalent} if $\pi_2(a)=U\pi_1(a)U\ad$ for a unitary $U:\HH_1\to\HH_2$. They are \emph{disjoint} if no non-zero subrepresentation of one is equivalent to a subrepresentation of the other, and \emph{quasi-equivalent} if each is equivalent to a multiple (direct sum of copies) of a subrepresentation of the other; this holds iff the map $\pi_1(a)\mapsto\pi_2(a)$ extends to an isomorphism of the von Neumann algebras $\pi_1(\A)''\cong\pi_2(\A)''$.

For a representation $\pi$ the \emph{folium} is the set of states given by density matrices in $\pi$ (normal states of $\pi(\A)''$). Two states $\omega_1,\omega_2$ have quasi-equivalent GNS representations iff they lie in the same folium. Physically:

\begin{center}
\emph{States in the same folium differ by ``finitely many'' local excitations: they are experimentally indistinguishable by local operations on a finite region, up to an arbitrarily small error, and they describe the same phase.}
\end{center}
States in different (disjoint) folia describe different phases or sectors: they cannot be obtained from each other by any local operations (or normal perturbations). For \emph{pure} states, quasi-equivalence of GNS representations is the same as unitary equivalence.

\section{A decisive tool: central macroscopic observables}

Recall from \cref{ch:topologies} that in the product representation $\pi_\theta$ the average magnetization $M_N=\frac1N\sum_{j\le N}\sigma^z_j$ converges strongly to the scalar $m(\theta)\id=\cos2\theta\,\id$. If $\pi_\theta$ and $\pi_{\theta'}$ were unitarily equivalent via $U$, then $U\big(m(\theta)\id\big)U\ad=\lim UM_NU\ad=\lim M_N=m(\theta')\id$, forcing $m(\theta)=m(\theta')$. Hence
\begin{equation}
\theta\ne\theta'\Longrightarrow\pi_\theta\not\simeq\pi_{\theta'},\quad\text{in fact they are disjoint.}
\end{equation}
\emph{General principle:} a strong limit of local observables which is a scalar in one representation and a different scalar in another proves inequivalence. In the weak closure $\pi(\A)''$, such limits are \emph{central} elements; they are the ``order parameters'' that label phases.

\section{A computable criterion: product states (Kakutani)}

Let $\omega=\bigotimes_j\omega_j$ and $\omega'=\bigotimes_j\omega'_j$ be infinite product \emph{vector} states of the spin chain, with local vectors $\psi_j,\psi_j'\in\C^2$.

\begin{theorem}[Kakutani]
$\pi_\omega$ and $\pi_{\omega'}$ are unitarily equivalent if $\sum_j\big(1-|\braket{\psi_j}{\psi_j'}|\big)<\infty$ (equivalently $\prod_j|\braket{\psi_j}{\psi'_j}|>0$), and disjoint otherwise.
\end{theorem}
For the states $\psi_\theta$ at every site, $|\braket{\psi_\theta}{\psi_{\theta'}}|=\cos(\theta-\theta')<1$ is constant, so the product $\to0$ and they are disjoint, consistent with the above. The criterion says: two states are equivalent iff they differ \emph{significantly} on only finitely many sites. This is the lattice version of ``a finite number of excitations'' versus ``a finite density of excitations''.

\begin{whatwrong}
A state which differs from another by a small amount at every site, however small, is in a different sector (for fixed non-summable perturbation). This is why thermodynamic limits cannot be treated as small perturbations of ground states, and why exponentiating a perturbation $\ee^{-\beta H_{\rm int}}$ in infinite volume may leave the Hilbert space.
\end{whatwrong}

\section{Spontaneous symmetry breaking}

Let $G$ act on $\A$ by automorphisms $\alpha_g$ (for example the spin flip $\sigma^z_j\mapsto-\sigma^z_j$, $\sigma^{x,y}_j\mapsto\sigma^{x,y}_j$ combined with a rotation, which maps $\psi_\theta\mapsto\psi_{\pi/2-\theta}$ and $m\to-m$). A state $\omega$ is \emph{invariant} if $\omega\circ\alpha_g=\omega$. If $\omega$ is invariant then $U(g)\pi_\omega(a)\Omega=\pi_\omega(\alpha_g(a))\Omega$ defines a unitary representation implementing $\alpha_g$ on $\HH_\omega$, leaving $\Omega$ fixed.

\begin{definition}
The symmetry $\alpha_g$ is \emph{spontaneously broken} in the representation $\pi_\omega$ if it is \emph{not} unitarily implemented there, i.e.\ if $\omega\circ\alpha_g$ is not in the folium of $\omega$.
\end{definition}

\emph{Example:} In the $m=+m_0$ ground state of an infinite ferromagnet the spin flip maps it to the $m=-m_0$ state, which lies in a different folium (just as for $\psi_\theta$ vs $\psi_{\pi/2-\theta}$): the flip has no implementer on $\HH_+$. The symmetry of the algebra is not a symmetry of the Hilbert space; the ``vacua'' are degenerate yet superselected: an observer cannot superpose them, since no local observable connects them (this is the origin of ``cat'' states being absent in infinite volume). Goldstone's theorem and its consequences are theorems about the (non-)implementability and the structure of charges: the algebraic charge $Q_\Lambda=\int_\Lambda j^0$ has no weak limit on $\HH_\omega$ if the symmetry is broken, yet it acts on local operators by $[Q_\Lambda,a]\to\delta(a)$ (the infinitesimal automorphism).

\begin{keyidea}
Inequivalent representations are not pathologies; they are phases. Superselection sectors are the labels of folia, order parameters are central elements, and spontaneous symmetry breaking means the automorphism group is not unitarily implemented.
\end{keyidea}

\section{Thermal vs vacuum states and an outlook}

For the Weyl algebra of a free field, the thermal representation at $T>0$ is inequivalent to the vacuum representation (\cref{ch:weyl}): the thermal state has a non-zero density of particles, so $\sum_k\abs{\beta_k}^2=\infty$ in infinite volume, and by Kakutani-type reasoning the two folia are disjoint. This is the algebraic fact behind ``the Unruh/thermal state is not a state in the vacuum Fock space'', and it will return as the statement that the modular flow of the Rindler wedge is not implemented by an element of the algebra (\cref{ch:rs_bw}). In algebraic QFT, superselection sectors for charges, for example localized charged states, are the content of the DHR analysis (\cref{ch:dhr}).

\section*{Exercises}
\begin{exercise}
Prove Kakutani's criterion in the easy direction: if $\sum_j(1-|\braket{\psi_j}{\psi_j'}|)<\infty$ show that $\ket{\Psi'}$ can be approximated by vectors obtained from $\ket{\Psi}$ by acting with local operators.
\end{exercise}
\begin{exercise}
Let $\omega_{p}$ be the product state with $\rho_j=\mathrm{diag}(p,1-p)$ at every site, $p\ne1/2$. Show that $\omega_p$ and $\omega_{1/2}=\tau$ are disjoint using $M_N$. What does this say about whether a state with a finite density of excitations can be reached from the tracial state by local operations?
\end{exercise}
\begin{exercise}
For the product state $\psi_\theta=\bigotimes_j(\cos\theta\ket\uparrow+\sin\theta\ket\downarrow)$, $\theta\in(0,\pi/4)$, show that $\psi_\theta$ is not invariant under the spin-flip automorphism $\sigma^z_j\mapsto-\sigma^z_j$ (with $\sigma^x_j\mapsto-\sigma^x_j$ so that the algebra relations hold). Show it is not unitarily implementable on $\HH_\theta$.
\end{exercise}
\begin{exercise}
Why can a finite system never exhibit spontaneous symmetry breaking in the sense above? (Use Stone--von Neumann or the fact that $\BH$ has only one irreducible representation when $\dim\HH<\infty$.)
\end{exercise}
